\medskip

Cet exercice contient 5 affirmations.

Pour chaque affirmation, répondre par VRAI ou FAUX en justifiant la réponse.

Toute absence de justification ou justification incorrecte ne sera pas prise en compte dans la notation.

\medskip

\textbf{Partie 1}

\medskip

On considère la suite $(u_n)$ définie par:

\begin{center}$u_0 = 10$\, et pour tout entier naturel $n,\,u_{n+1} = \dfrac13 u_n + 2$.\end{center}

\medskip

\begin{enumerate}
\item \textbf{Affirmation 1} : La suite $(u_n)$ est décroissante minorée par 0.
\item \textbf{Affirmation 2} : $\displaystyle\lim_{n \to + \infty} u_n = 0$.
\item \textbf{Affirmation 3} : La suite $(v_n)$ définie pour tout entier naturel $n$ par $v_n = u_n - 3$ est géométrique.
\end{enumerate}

\medskip

\textbf{Partie 2}

\medskip

On considère l'équation différentielle $(E) :\, y' = \dfrac32 y + 2$ d'inconnue $y$, fonction définie et 
dérivable sur $\R$

\medskip

\begin{enumerate}
\item \textbf{Affirmation 4} : Il existe une fonction constante solution de l'équation différentielle (E).
\item Dans un repère orthonormé \Oij{} on note $\mathcal{C}_f$ la courbe représentative de la fonction $f$ solution de $(E)$ telle  que $f(0) = 0$.

\item \textbf{Affirmation 5} : La tangente au point d'abscisse 1 de $\mathcal{C}_f$ a pour coefficient directeur $2 \e^{\frac32}$.
\end{enumerate}

\bigskip
